% ============================================================================
% CAPÍTULO 15: EMERGING TOPICS
% Bioinformática para Biologia de Sistemas
% ============================================================================

% ============================================================================
% TEMPLATE BASE PARA APRESENTAÇÕES EM BEAMER
% Bioinformática para Biologia de Sistemas
% ============================================================================

\documentclass[12pt, aspectratio=169]{beamer}

% ============================================================================
% PACOTES ESSENCIAIS
% ============================================================================
\usepackage[utf8]{inputenc}
\usepackage[T1]{fontenc}
\usepackage[brazilian]{babel}
\usepackage{amsmath, amssymb, amsthm}
\usepackage{graphicx}
\usepackage{xcolor}
\usepackage{tikz}
\usetikzlibrary{arrows.meta, shapes, positioning, calc, decorations.pathreplacing, decorations.shapes, decorations.pathmorphing}
\usepackage{listings}
\usepackage{booktabs}
\usepackage{multirow}
\usepackage{hyperref}
\usepackage{chemfig}
\usepackage{chemarr}

% ============================================================================
% CONFIGURAÇÃO DO TEMA
% ============================================================================
\usetheme{Madrid}
\usecolortheme{default}

% Cores personalizadas
\definecolor{azulescuro}{RGB}{0, 51, 102}
\definecolor{azulclaro}{RGB}{51, 153, 255}
\definecolor{verdeacido}{RGB}{102, 204, 0}
\definecolor{laranjacel}{RGB}{255, 153, 51}
\definecolor{roxodna}{RGB}{153, 51, 255}

\setbeamercolor{structure}{fg=azulescuro}
\setbeamercolor{palette primary}{bg=azulescuro,fg=white}
\setbeamercolor{palette secondary}{bg=azulclaro,fg=white}
\setbeamercolor{palette tertiary}{bg=azulescuro,fg=white}
\setbeamercolor{palette quaternary}{bg=azulclaro,fg=white}
\setbeamercolor{block title}{bg=azulescuro,fg=white}
\setbeamercolor{block body}{bg=azulclaro!10,fg=black}
\setbeamercolor{block title alerted}{bg=red!80,fg=white}
\setbeamercolor{block body alerted}{bg=red!10,fg=black}
\setbeamercolor{block title example}{bg=verdeacido!80,fg=white}
\setbeamercolor{block body example}{bg=verdeacido!10,fg=black}

% ============================================================================
% CONFIGURAÇÃO DE CÓDIGO
% ============================================================================
\lstset{
    language=Python,
    basicstyle=\ttfamily\small,
    keywordstyle=\color{azulescuro}\bfseries,
    commentstyle=\color{verdeacido},
    stringstyle=\color{laranjacel},
    numbers=left,
    numberstyle=\tiny\color{gray},
    stepnumber=1,
    numbersep=5pt,
    backgroundcolor=\color{white},
    showspaces=false,
    showstringspaces=false,
    showtabs=false,
    frame=single,
    rulecolor=\color{black},
    tabsize=2,
    captionpos=b,
    breaklines=true,
    breakatwhitespace=false,
    escapeinside={\%*}{*)},
    xleftmargin=10pt,
    xrightmargin=10pt
}

% Definir estilo para R
\lstdefinestyle{Rstyle}{
    language=R,
    basicstyle=\ttfamily\small,
    keywordstyle=\color{azulescuro}\bfseries,
    commentstyle=\color{verdeacido},
    stringstyle=\color{laranjacel}
}

% Definir estilo para MATLAB
\lstdefinestyle{MATLABstyle}{
    language=Matlab,
    basicstyle=\ttfamily\small,
    keywordstyle=\color{azulescuro}\bfseries,
    commentstyle=\color{verdeacido},
    stringstyle=\color{laranjacel}
}

% ============================================================================
% COMANDOS PERSONALIZADOS
% ============================================================================

% Comando para equações destacadas
\newcommand{\destaque}[1]{%
    \begin{alertblock}{}
        \begin{center}
            #1
        \end{center}
    \end{alertblock}
}

% Comando para definições
\newcommand{\definicao}[2]{%
    \begin{block}{Definição: #1}
        #2
    \end{block}
}

% Comando para exemplos
\newcommand{\exemplo}[2]{%
    \begin{exampleblock}{Exemplo: #1}
        #2
    \end{exampleblock}
}

% Comando para observações importantes
\newcommand{\importante}[1]{%
    \begin{alertblock}{Importante!}
        #1
    \end{alertblock}
}

% Comandos matemáticos úteis
\newcommand{\R}{\mathbb{R}}
\newcommand{\N}{\mathbb{N}}
\newcommand{\Z}{\mathbb{Z}}
\newcommand{\dif}{\mathrm{d}}
\newcommand{\parcial}[2]{\frac{\partial #1}{\partial #2}}
\newcommand{\derivada}[2]{\frac{\dif #1}{\dif #2}}
\newcommand{\vect}[1]{\boldsymbol{#1}}

% Comandos biológicos
\newcommand{\gene}[1]{\textit{#1}}
\newcommand{\proteina}[1]{\textsc{#1}}
\newcommand{\especie}[1]{\textit{#1}}

% ============================================================================
% CONFIGURAÇÃO DE RODAPÉ
% ============================================================================
\setbeamertemplate{footline}{
  \leavevmode%
  \hbox{%
  \begin{beamercolorbox}[wd=.33\paperwidth,ht=2.25ex,dp=1ex,center]{author in head/foot}%
    \usebeamerfont{author in head/foot}\insertshortauthor
  \end{beamercolorbox}%
  \begin{beamercolorbox}[wd=.34\paperwidth,ht=2.25ex,dp=1ex,center]{title in head/foot}%
    \usebeamerfont{title in head/foot}\insertshorttitle
  \end{beamercolorbox}%
  \begin{beamercolorbox}[wd=.33\paperwidth,ht=2.25ex,dp=1ex,right]{date in head/foot}%
    \usebeamerfont{date in head/foot}\insertshortdate{}\hspace*{2em}
    \insertframenumber{} / \inserttotalframenumber\hspace*{2ex}
  \end{beamercolorbox}}%
  \vskip0pt%
}

% ============================================================================
% CONFIGURAÇÃO DE NAVEGAÇÃO
% ============================================================================
\setbeamertemplate{navigation symbols}{}

% ============================================================================
% CONFIGURAÇÃO DE ITEMIZE/ENUMERATE
% ============================================================================
\setbeamertemplate{itemize items}[circle]
\setbeamertemplate{enumerate items}[default]

% ============================================================================
% FIM DO TEMPLATE
% ============================================================================


% ============================================================================
% INFORMAÇÕES DO DOCUMENTO
% ============================================================================
\title[Cap. 15: Emerging Topics]{Capítulo 15: Emerging Topics}
\subtitle{Bioinformática para Biologia de Sistemas}
\author{Seu Nome}
\institute{Sua Instituição}
\date{\today}

\begin{document}

% ============================================================================
% SLIDE DE TÍTULO
% ============================================================================
\begin{frame}
\titlepage
\end{frame}

% ============================================================================
% SUMÁRIO
% ============================================================================
\begin{frame}{Sumário}
\tableofcontents
\end{frame}

% ============================================================================
% SEÇÃO 1: INTRODUÇÃO
% ============================================================================
\section{Introdução}

\begin{frame}{Visão Geral do Capítulo}
\begin{block}{Objetivos de Aprendizagem}
\begin{itemize}
    \item Explorar tópicos emergentes em biologia de sistemas
    \item Compreender modelagem multi-escala
    \item Conhecer análise de célula única e transcriptômica espacial
    \item Aprender sobre inteligência artificial aplicada à biologia
    \item Discutir doenças complexas, envelhecimento e medicina de precisão
    \item Vislumbrar direções futuras da área
\end{itemize}
\end{block}

\vspace{0.3cm}

\importante{
Este é o capítulo final do curso, onde integramos conhecimentos anteriores e olhamos para o futuro da biologia de sistemas!
}
\end{frame}

\begin{frame}{Evolução Rápida da Biologia de Sistemas}
\begin{center}
\begin{tikzpicture}[node distance=2.5cm, auto,
    box/.style={rectangle, draw, text width=2cm, text centered, rounded corners, minimum height=1cm, font=\small}]

    \node[box, fill=roxodna!30] (genomics) {Genômica\\(2000s)};
    \node[box, fill=laranjacel!30, right of=genomics] (omics) {Multi-\\Omics\\(2010s)};
    \node[box, fill=verdeacido!30, right of=omics] (sc) {Single-\\Cell\\(2015+)};
    \node[box, fill=azulclaro!30, right of=sc] (ai) {AI +\\Systems\\(2020+)};

    \draw[->, thick] (genomics) -- (omics);
    \draw[->, thick] (omics) -- (sc);
    \draw[->, thick] (sc) -- (ai);
\end{tikzpicture}
\end{center}

\vspace{0.3cm}

\begin{exampleblock}{Tendências Atuais}
\begin{itemize}
    \item Integração de múltiplas camadas ômicas
    \item Resolução de célula única e espacial
    \item Modelagem multi-escala (molécula $\rightarrow$ organismo)
    \item IA para descoberta e predição
    \item Medicina personalizada e digital twins
\end{itemize}
\end{exampleblock}
\end{frame}

% ============================================================================
% SEÇÃO 2: MODELAGEM MULTI-ESCALA
% ============================================================================
\section{Modelagem Multi-Escala}

\begin{frame}{15.1 Modelagem Multi-Escala}
\definicao{Modelagem Multi-Escala}{
Integração de modelos computacionais em diferentes níveis de organização biológica: molecular, celular, tecidual, órgão e organismo completo.
}

\vspace{0.3cm}

\begin{block}{Desafios da Integração Multi-Escala}
\begin{enumerate}
    \item \textbf{Escalas temporais:} ms (cinética enzimática) $\rightarrow$ anos (envelhecimento)
    \item \textbf{Escalas espaciais:} nm (proteínas) $\rightarrow$ m (organismo)
    \item \textbf{Complexidade:} $10^{12}$ moléculas $\rightarrow$ trilhões de células
    \item \textbf{Heterogeneidade:} variabilidade celular e individual
    \item \textbf{Computação:} custo proibitivo de simulações completas
\end{enumerate}
\end{block}

\importante{
Nenhum modelo único pode capturar todos os níveis simultaneamente!
}
\end{frame}

\begin{frame}{Hierarquia Multi-Escala}
\begin{center}
\begin{tikzpicture}[scale=0.9]
    % Multi-scale hierarchy
    \foreach \y/\scale/\color/\time in {
        0/Organismo/azulclaro/anos,
        1.2/Órgão/verdeacido/meses,
        2.4/Tecido/laranjacel/dias,
        3.6/Célula/roxodna/horas,
        4.8/Via/azulescuro/minutos,
        6/Molécula/red/ms-s
    } {
        \draw[fill=\color!30, draw=\color!80, thick]
            (-3+\y*0.3,-\y*0.8) rectangle (3-\y*0.3,-\y*0.8-0.6);
        \node[font=\small] at (0,-\y*0.8-0.3) {\scale};
        \node[right, font=\tiny] at (3.2-\y*0.3,-\y*0.8-0.3) {\time};
    }

    % Arrows
    \foreach \y in {0,1.2,2.4,3.6,4.8} {
        \draw[->, thick, gray] (-3.5,-\y*0.8-0.7) -- (-3.5,-\y*0.8-1.1);
    }

    \node[left, font=\small, rotate=90] at (-4,-2.4) {Crescente complexidade};
\end{tikzpicture}
\end{center}

\begin{exampleblock}{Exemplo: Metabolismo da Glicose}
\begin{itemize}
    \item \textbf{Molecular:} cinética de hexoquinase
    \item \textbf{Via:} glicólise completa
    \item \textbf{Celular:} hepatócito produzindo ATP
    \item \textbf{Órgão:} fígado regulando glicemia
    \item \textbf{Organismo:} homeostase sistêmica
\end{itemize}
\end{exampleblock}
\end{frame}

\begin{frame}{Modelos Baseados em Agentes (ABM)}
\definicao{Agent-Based Model}{
Modelo computacional onde entidades autônomas (agentes) seguem regras simples, e comportamentos complexos emergem das interações.
}

\vspace{0.3cm}

\begin{columns}
\column{0.5\textwidth}
\textbf{Características:}
\begin{itemize}
    \item Agentes autônomos
    \item Regras locais simples
    \item Interações entre agentes
    \item Ambiente compartilhado
    \item Emergência de padrões
\end{itemize}

\column{0.5\textwidth}
\textbf{Aplicações:}
\begin{itemize}
    \item Crescimento tumoral
    \item Resposta imune
    \item Morfogênese
    \item Dinâmica de populações
    \item Epidemiologia
\end{itemize}
\end{columns}

\vspace{0.3cm}

\begin{exampleblock}{Exemplo: Tumor}
Cada célula é um agente que pode: proliferar, migrar, morrer, ou entrar em quiescência baseado em sinais locais (nutrientes, fatores de crescimento).
\end{exampleblock}
\end{frame}

\begin{frame}[fragile]{Exemplo Simples de ABM em Python}
\begin{lstlisting}[language=Python, caption=Agent-Based Model básico, basicstyle=\ttfamily\footnotesize]
import numpy as np
import matplotlib.pyplot as plt

class Cell:
    """Agente representando uma celula"""
    def __init__(self, x, y, cell_type='normal'):
        self.x = x
        self.y = y
        self.type = cell_type
        self.alive = True

    def proliferate(self, environment):
        """Celula pode proliferar se houver espaco"""
        if environment.nutrient_level[self.x, self.y] > 0.5:
            # Criar celula filha adjacente
            new_x = self.x + np.random.randint(-1, 2)
            new_y = self.y + np.random.randint(-1, 2)
            if environment.is_empty(new_x, new_y):
                return Cell(new_x, new_y, self.type)
        return None

    def move(self, environment):
        """Celula se move aleatoriamente"""
        dx, dy = np.random.randint(-1, 2, 2)
        new_x, new_y = self.x + dx, self.y + dy
        if environment.is_valid(new_x, new_y):
            self.x, self.y = new_x, new_y
\end{lstlisting}
\end{frame}

\begin{frame}{Modelos Híbridos: Discreto + Contínuo}
\begin{block}{Abordagem Híbrida}
Combina modelos discretos (células individuais) com modelos contínuos (difusão de nutrientes, fatores de crescimento).
\end{block}

\begin{center}
\begin{tikzpicture}[scale=0.8]
    % Continuous field (gradient)
    \foreach \x in {0,...,5} {
        \foreach \y in {0,...,5} {
            \pgfmathsetmacro{\conc}{100 - 10*sqrt((\x-2.5)^2 + (\y-2.5)^2)}
            \definecolor{fieldcolor}{rgb}{1,\conc/100,\conc/100}
            \fill[fieldcolor] (\x*0.8,\y*0.8) rectangle (\x*0.8+0.8,\y*0.8+0.8);
        }
    }

    % Discrete agents (cells)
    \foreach \pos in {(1.2,1.2), (2.4,1.6), (1.8,2.8), (3.2,2.4), (2.0,3.6)} {
        \fill[azulescuro] \pos circle (0.2);
    }

    \node[below, font=\small] at (2,-0.5) {Campo contínuo (nutrientes)};
    \node[above, font=\small] at (2,5.5) {Agentes discretos (células)};
\end{tikzpicture}
\end{center}

\textbf{Equações acopladas:}
\begin{itemize}
    \item \textbf{Contínuo:} $\parcial{C}{t} = D\nabla^2 C - \sum_i \text{uptake}_i(C)$
    \item \textbf{Discreto:} crescimento celular baseado em $C(x_i, y_i)$
\end{itemize}
\end{frame}

\begin{frame}{Órgãos Virtuais e Fisiologia de Corpo Inteiro}
\begin{columns}
\column{0.5\textwidth}
\textbf{Órgãos Virtuais:}
\begin{itemize}
    \item Coração virtual (eletrof.)
    \item Fígado virtual (metabolismo)
    \item Pulmão virtual (ventilação)
    \item Rim virtual (filtração)
\end{itemize}

\vspace{0.3cm}

\textbf{Projetos Internacionais:}
\begin{itemize}
    \item \textbf{Virtual Physiological Human}
    \item \textbf{Physiome Project}
    \item \textbf{Human Brain Project}
\end{itemize}

\column{0.5\textwidth}
\begin{center}
\begin{tikzpicture}[scale=0.7]
    % Body outline
    \draw[thick, azulescuro] (0,0) ellipse (1 and 2);

    % Organs
    \fill[red!50] (0,1) circle (0.3); % brain
    \node[font=\tiny] at (0,1) {Cérebro};

    \fill[laranjacel!50] (-0.5,0.2) circle (0.2); % heart
    \node[font=\tiny] at (-0.5,0.2) {C};

    \fill[verdeacido!50] (0,-0.3) circle (0.25); % liver
    \node[font=\tiny] at (0,-0.3) {Fíg};

    \fill[azulclaro!50] (0.5,-0.8) circle (0.15); % kidney
    \fill[azulclaro!50] (-0.5,-0.8) circle (0.15);
    \node[font=\tiny] at (0,-1.2) {Rins};

    % Connections
    \draw[<->, thick, red] (-0.5,0.2) -- (0,1);
    \draw[<->, thick, red] (0,-0.3) -- (-0.5,0.2);
\end{tikzpicture}
\end{center}

\begin{alertblock}{Desafio}
Integrar dinâmicas de múltiplos órgãos acoplados!
\end{alertblock}
\end{columns}
\end{frame}

\begin{frame}{Exemplo: Metabolismo Hepático Multi-Escala}
\begin{block}{Níveis Integrados}
\begin{enumerate}
    \item \textbf{Molecular:} cinética de enzimas (CYP450, glicoquinase)
    \item \textbf{Celular:} hepatócitos individuais (zonas periportal/pericentral)
    \item \textbf{Lobular:} gradiente de O$_2$ e nutrientes
    \item \textbf{Órgão:} fluxo sanguíneo, metabolismo sistêmico
    \item \textbf{Organismo:} regulação hormonal (insulina, glucagon)
\end{enumerate}
\end{block}

\vspace{0.3cm}

\begin{exampleblock}{Modelagem}
\begin{itemize}
    \item ODEs para metabolismo celular
    \item PDEs para difusão espacial
    \item ABM para heterogeneidade hepatócito
    \item Controle sistêmico (feedback hormonal)
\end{itemize}
\end{exampleblock}

\importante{
Permite simular resposta a drogas, doenças hepáticas, cirrose!
}
\end{frame}

% ============================================================================
% SEÇÃO 3: BIOLOGIA DE SISTEMAS DE CÉLULA ÚNICA
% ============================================================================
\section{Biologia de Sistemas de Célula Única}

\begin{frame}{15.2 Biologia de Sistemas de Célula Única}
\begin{block}{Revolução da Célula Única}
Tecnologias que permitem medir milhares de moléculas em células individuais, revelando heterogeneidade antes invisível.
\end{block}

\begin{columns}
\column{0.5\textwidth}
\textbf{Técnicas Principais:}
\begin{itemize}
    \item scRNA-seq
    \item scATAC-seq (cromatina)
    \item scProteomics (CyTOF)
    \item scMetabolomics
    \item Multi-omic single-cell
\end{itemize}

\column{0.5\textwidth}
\textbf{Aplicações:}
\begin{itemize}
    \item Descoberta de tipos celulares
    \item Trajetórias de diferenciação
    \item Heterogeneidade tumoral
    \item Desenvolvimento embrionário
    \item Resposta imune
\end{itemize}
\end{columns}

\vspace{0.3cm}

\importante{
A célula única é a unidade fundamental da vida. Tecidos aparentemente homogêneos são, na verdade, mosaicos heterogêneos!
}
\end{frame}

\begin{frame}{Single-Cell RNA-seq (scRNA-seq)}
\begin{center}
\begin{tikzpicture}[node distance=1.8cm, auto,
    box/.style={rectangle, draw, text width=2.5cm, text centered, rounded corners, minimum height=0.8cm, font=\small}]

    \node[box, fill=roxodna!30] (tissue) {Tecido\\Heterogêneo};
    \node[box, fill=laranjacel!30, below of=tissue] (dissoc) {Dissociação\\em Células};
    \node[box, fill=verdeacido!30, below of=dissoc] (capture) {Captura de\\Células};
    \node[box, fill=azulclaro!30, below of=capture] (seq) {Sequenciamento\\cDNA};
    \node[box, fill=yellow!30, below of=seq] (analysis) {Análise\\Bioinformática};

    \draw[->, thick] (tissue) -- (dissoc);
    \draw[->, thick] (dissoc) -- (capture);
    \draw[->, thick] (capture) -- (seq);
    \draw[->, thick] (seq) -- (analysis);

    % Annotations
    \node[right, font=\tiny, text width=2.5cm] at (3.5,-1.8) {Droplet-based\\(10x Genomics)};
    \node[right, font=\tiny, text width=2.5cm] at (3.5,-7.2) {Clustering, DE,\\Trajectories};
\end{tikzpicture}
\end{center}
\end{frame}

\begin{frame}{Identificação de Tipos Celulares e Clustering}
\begin{block}{Pipeline de Análise scRNA-seq}
\begin{enumerate}
    \item \textbf{Controle de qualidade:} remover células de baixa qualidade
    \item \textbf{Normalização:} corrigir profundidade de sequenciamento
    \item \textbf{Seleção de genes:} genes altamente variáveis
    \item \textbf{Redução de dimensionalidade:} PCA, UMAP, t-SNE
    \item \textbf{Clustering:} Louvain, Leiden
    \item \textbf{Anotação:} marcadores conhecidos
\end{enumerate}
\end{block}

\vspace{0.3cm}

\begin{exampleblock}{Ferramentas}
\begin{itemize}
    \item \textbf{Scanpy} (Python)
    \item \textbf{Seurat} (R)
    \item \textbf{Monocle} (trajectories)
    \item \textbf{Cell Ranger} (10x)
\end{itemize}
\end{exampleblock}
\end{frame}

\begin{frame}[fragile]{Análise scRNA-seq com Scanpy}
\begin{lstlisting}[language=Python, caption=Análise básica com scanpy, basicstyle=\ttfamily\footnotesize]
import scanpy as sc
import numpy as np

# Carregar dados
adata = sc.read_10x_mtx('filtered_feature_bc_matrix/')

# QC: filtrar celulas e genes
sc.pp.filter_cells(adata, min_genes=200)
sc.pp.filter_genes(adata, min_cells=3)

# Calcular % mitocondriais
adata.var['mt'] = adata.var_names.str.startswith('MT-')
sc.pp.calculate_qc_metrics(
    adata, qc_vars=['mt'], inplace=True
)

# Filtrar celulas ruins
adata = adata[adata.obs.pct_counts_mt < 5, :]

# Normalizar e log-transformar
sc.pp.normalize_total(adata, target_sum=1e4)
sc.pp.log1p(adata)

# Identificar genes altamente variaveis
sc.pp.highly_variable_genes(adata, n_top_genes=2000)

# PCA
sc.tl.pca(adata, svd_solver='arpack')
\end{lstlisting}
\end{frame}

\begin{frame}[fragile]{Clustering e Visualização}
\begin{lstlisting}[language=Python, caption=Clustering e UMAP, basicstyle=\ttfamily\footnotesize]
# Construir grafo de vizinhos
sc.pp.neighbors(adata, n_neighbors=10, n_pcs=40)

# Clustering com algoritmo Leiden
sc.tl.leiden(adata, resolution=0.5)

# UMAP para visualizacao
sc.tl.umap(adata)

# Plotar
sc.pl.umap(adata, color=['leiden', 'n_genes'])

# Encontrar genes marcadores
sc.tl.rank_genes_groups(adata, 'leiden', method='wilcoxon')
sc.pl.rank_genes_groups(adata, n_genes=20, sharey=False)

# Anotar clusters manualmente
cluster_names = {
    '0': 'T cells',
    '1': 'B cells',
    '2': 'Monocytes',
    '3': 'NK cells',
    '4': 'Dendritic cells'
}
adata.obs['cell_type'] = adata.obs['leiden'].map(cluster_names)
sc.pl.umap(adata, color='cell_type')
\end{lstlisting}
\end{frame}

\begin{frame}{Pseudotempo e Inferência de Trajetórias}
\definicao{Pseudotempo}{
Ordem inferida de células ao longo de um processo dinâmico (diferenciação, ativação) baseada em similaridade transcricional.
}

\vspace{0.3cm}

\begin{center}
\begin{tikzpicture}[scale=0.9]
    % Start cell
    \fill[verdeacido!50] (0,0) circle (0.3);
    \node[below] at (0,-0.5) {Stem};

    % Trajectory 1
    \foreach \x in {1,...,4} {
        \fill[verdeacido!50, opacity={1-\x*0.2}] (\x*1.5,1) circle (0.2);
        \draw[->, thick] (\x*1.5-0.7,0.9) -- (\x*1.5-0.3,1);
    }
    \node[above] at (6,1) {Tipo A};

    % Trajectory 2
    \foreach \x in {1,...,4} {
        \fill[azulclaro!50, opacity={1-\x*0.2}] (\x*1.5,-1) circle (0.2);
        \draw[->, thick] (\x*1.5-0.7,-0.9) -- (\x*1.5-0.3,-1);
    }
    \node[below] at (6,-1) {Tipo B};

    % Connection from start
    \draw[->, thick] (0.3,0) -- (0.8,1);
    \draw[->, thick] (0.3,0) -- (0.8,-1);

    % Pseudotime axis
    \draw[->, thick] (-0.5,-2.5) -- (6.5,-2.5);
    \node[below] at (3,-2.5) {Pseudotempo $\rightarrow$};
\end{tikzpicture}
\end{center}

\textbf{Algoritmos:} Monocle, Slingshot, PAGA, RNA velocity
\end{frame}

\begin{frame}{Transcriptômica Espacial}
\begin{block}{Resolvendo a Localização}
scRNA-seq perde informação espacial. Novas tecnologias preservam contexto espacial:
\end{block}

\begin{columns}
\column{0.5\textwidth}
\textbf{Tecnologias:}
\begin{itemize}
    \item \textbf{Visium} (10x): arrays espaciais
    \item \textbf{MERFISH}: imaging baseado em FISH
    \item \textbf{seqFISH+}: multiplexado
    \item \textbf{Slide-seq}: beads espaciais
    \item \textbf{Stereo-seq}: ultra-alta resolução
\end{itemize}

\column{0.5\textwidth}
\textbf{Aplicações:}
\begin{itemize}
    \item Arquitetura tecidual
    \item Microambiente tumoral
    \item Organização cerebral
    \item Nicho de células-tronco
    \item Desenvolvimento embrionário
\end{itemize}
\end{columns}

\vspace{0.3cm}

\importante{
Combina melhor dos dois mundos: expressão gênica detalhada + localização espacial!
}
\end{frame>

\begin{frame}{Comunicação Célula-Célula}
\begin{block}{Inferindo Interações}
Análise de pares ligante-receptor para identificar comunicação entre tipos celulares.
\end{block}

\begin{center}
\begin{tikzpicture}[scale=0.8]
    % Cell types
    \node[draw, circle, fill=verdeacido!30, minimum size=1.5cm] (cellA) at (0,0) {Célula A};
    \node[draw, circle, fill=azulclaro!30, minimum size=1.5cm] (cellB) at (4,0) {Célula B};

    % Ligand-receptor pairs
    \draw[->, thick, laranjacel, line width=2pt] (0.75,0.3) -- node[above, font=\small] {VEGF} (3.25,0.3);
    \node[right, font=\tiny] at (4.75,0.3) {VEGFR};

    \draw[->, thick, roxodna, line width=2pt] (3.25,-0.3) -- node[below, font=\small] {PD-L1} (0.75,-0.3);
    \node[left, font=\tiny] at (-0.75,-0.3) {PD-1};
\end{tikzpicture}
\end{center>

\vspace{0.3cm}

\begin{exampleblock}{Ferramentas}
\begin{itemize}
    \item \textbf{CellPhoneDB}
    \item \textbf{NicheNet}
    \item \textbf{CellChat}
    \item \textbf{LIANA}
\end{itemize}
\end{exampleblock}

\textbf{Aplicação:} microambiente tumoral, resposta imune, nicho de células-tronco
\end{frame}

% ============================================================================
% SEÇÃO 4: INTELIGÊNCIA ARTIFICIAL EM BIOLOGIA DE SISTEMAS
% ============================================================================
\section{Inteligência Artificial em Biologia de Sistemas}

\begin{frame}{15.3 IA em Biologia de Sistemas}
\begin{block}{A Revolução da IA}
Machine learning e deep learning estão transformando a biologia de sistemas, permitindo análise de dados massivos e descoberta de padrões complexos.
\end{block>

\begin{columns}
\column{0.5\textwidth}
\textbf{Machine Learning:}
\begin{itemize}
    \item Supervisionado (classificação)
    \item Não-supervisionado (clustering)
    \item Reforço (otimização)
    \item Ensemble methods
\end{itemize}

\column{0.5\textwidth}
\textbf{Deep Learning:}
\begin{itemize}
    \item Redes neurais profundas
    \item CNNs (imagens)
    \item RNNs/Transformers (sequências)
    \item GNNs (grafos/redes)
    \item GANs (geração)
\end{itemize}
\end{columns}

\vspace{0.3cm}

\importante{
IA não substitui modelagem mecanística, mas complementa: híbrido é o futuro!
}
\end{frame}

\begin{frame}{Aprendizado Profundo para Dados Biológicos}
\begin{center}
\begin{tikzpicture}[scale=0.8]
    % Input layer
    \foreach \y in {1,...,4} {
        \node[draw, circle, fill=roxodna!30, minimum size=0.6cm] (i\y) at (0,\y) {};
    }
    \node[below] at (0,0.5) {Input};

    % Hidden layer 1
    \foreach \y in {0.5,...,4.5} {
        \node[draw, circle, fill=laranjacel!30, minimum size=0.5cm] (h1\y) at (2,\y) {};
    }
    \node[below] at (2,0) {Hidden 1};

    % Hidden layer 2
    \foreach \y in {0.5,...,4.5} {
        \node[draw, circle, fill=laranjacel!30, minimum size=0.5cm] (h2\y) at (4,\y) {};
    }
    \node[below] at (4,0) {Hidden 2};

    % Output layer
    \foreach \y in {1.5,3.5} {
        \node[draw, circle, fill=verdeacido!30, minimum size=0.6cm] (o\y) at (6,\y) {};
    }
    \node[below] at (6,0) {Output};

    % Connections (simplified)
    \foreach \i in {1,...,4} {
        \foreach \h in {0.5,1.5,2.5,3.5,4.5} {
            \draw[->, thin, gray] (i\i) -- (h1\h);
        }
    }
\end{tikzpicture}
\end{center}

\textbf{Aplicações:}
\begin{itemize}
    \item Predição de expressão gênica
    \item Classificação de tipos celulares
    \item Predição de estrutura proteica (AlphaFold)
    \item Design de drogas
    \item Diagnóstico de imagens médicas
\end{itemize}
\end{frame}

\begin{frame}[fragile]{Rede Neural para Classificação}
\begin{lstlisting}[language=Python, caption=RNA simples com PyTorch, basicstyle=\ttfamily\footnotesize]
import torch
import torch.nn as nn
import torch.optim as optim

class GeneExpressionClassifier(nn.Module):
    """Rede neural para classificar amostras"""
    def __init__(self, input_dim, hidden_dim, output_dim):
        super().__init__()
        self.fc1 = nn.Linear(input_dim, hidden_dim)
        self.relu = nn.ReLU()
        self.dropout = nn.Dropout(0.3)
        self.fc2 = nn.Linear(hidden_dim, hidden_dim // 2)
        self.fc3 = nn.Linear(hidden_dim // 2, output_dim)
        self.softmax = nn.Softmax(dim=1)

    def forward(self, x):
        x = self.fc1(x)
        x = self.relu(x)
        x = self.dropout(x)
        x = self.fc2(x)
        x = self.relu(x)
        x = self.fc3(x)
        return self.softmax(x)

# Inicializar modelo
model = GeneExpressionClassifier(
    input_dim=20000,  # genes
    hidden_dim=256,
    output_dim=5      # classes
)
\end{lstlisting}
\end{frame}

\begin{frame}{Graph Neural Networks (GNNs)}
\definicao{GNN}{
Redes neurais que operam diretamente em grafos, aprendendo representações de nós e arestas. Ideais para redes biológicas!
}

\vspace{0.3cm}

\begin{columns}
\column{0.5\textwidth}
\textbf{Aplicações:}
\begin{itemize}
    \item Predição de função de proteínas
    \item Descoberta de drogas
    \item Propriedades moleculares
    \item Redes regulatórias
    \item Interações proteína-proteína
\end{itemize}

\column{0.5\textwidth}
\begin{center}
\begin{tikzpicture}[scale=0.7]
    % Graph
    \foreach \i/\x/\y in {1/0/0, 2/2/1, 3/2/-1, 4/4/0} {
        \node[draw, circle, fill=azulclaro!30, minimum size=0.8cm] (n\i) at (\x,\y) {\i};
    }

    \draw[thick] (n1) -- (n2);
    \draw[thick] (n1) -- (n3);
    \draw[thick] (n2) -- (n4);
    \draw[thick] (n3) -- (n4);

    \node[below] at (2,-2) {Rede PPI};

    % Embedding
    \draw[->, ultra thick] (5,0) -- (6,0);
    \node[above] at (5.5,0) {GNN};

    % Embeddings
    \foreach \i/\y in {1/1.5, 2/0.5, 3/-0.5, 4/-1.5} {
        \fill[verdeacido!50] (7,\y) rectangle (8.5,\y+0.3);
    }
    \node[below] at (7.75,-2) {Embeddings};
\end{tikzpicture}
\end{center}
\end{columns>

\vspace{0.3cm}

\textbf{Modelos:} GCN, GraphSAGE, GAT, Message Passing NN
\end{frame}

\begin{frame}[fragile]{GNN Básico para Rede Molecular}
\begin{lstlisting}[language=Python, caption=GNN com PyTorch Geometric, basicstyle=\ttfamily\footnotesize]
import torch
from torch_geometric.nn import GCNConv, global_mean_pool

class MolecularGNN(torch.nn.Module):
    """GNN para predicao de propriedades moleculares"""
    def __init__(self, num_features, hidden_dim, num_classes):
        super().__init__()
        self.conv1 = GCNConv(num_features, hidden_dim)
        self.conv2 = GCNConv(hidden_dim, hidden_dim)
        self.conv3 = GCNConv(hidden_dim, hidden_dim)
        self.fc = torch.nn.Linear(hidden_dim, num_classes)

    def forward(self, x, edge_index, batch):
        # Node embeddings
        x = self.conv1(x, edge_index)
        x = torch.relu(x)
        x = self.conv2(x, edge_index)
        x = torch.relu(x)
        x = self.conv3(x, edge_index)

        # Graph-level pooling
        x = global_mean_pool(x, batch)

        # Prediction
        return self.fc(x)
\end{lstlisting}
\end{frame}

\begin{frame}{Transformers e AlphaFold}
\begin{block}{Mecanismo de Atenção}
Transformers usam "attention" para capturar dependências de longo alcance em sequências.
\end{block}

\begin{exampleblock}{AlphaFold 2 (DeepMind, 2020)}
\begin{itemize}
    \item Revolucionou predição de estrutura proteica
    \item Precisão próxima a estruturas experimentais
    \item Baseia-se em: transformers + MSA + embeddings estruturais
    \item Predisse estrutura de \textgreater200M proteínas (AlphaFold DB)
\end{itemize}
\end{exampleblock}

\vspace{0.3cm}

\importante{
AlphaFold é um dos maiores sucessos de IA em biologia!
}

\textbf{Sequência $\rightarrow$ MSA $\rightarrow$ Transformer $\rightarrow$ Estrutura 3D}
\end{frame}

\begin{frame}{Modelos Generativos: Design de Moléculas}
\begin{columns}
\column{0.5\textwidth}
\textbf{Generative Models:}
\begin{itemize}
    \item \textbf{VAE:} Variational Autoencoder
    \item \textbf{GAN:} Generative Adversarial Network
    \item \textbf{Diffusion:} Diffusion models
    \item \textbf{RL:} Reinforcement learning
\end{itemize>

\vspace{0.3cm}

\textbf{Aplicações:}
\begin{itemize}
    \item Design de proteínas
    \item Design de drogas (SMILES)
    \item Otimização de anticorpos
\end{itemize}

\column{0.5\textwidth}
\begin{center}
\begin{tikzpicture}[scale=0.7]
    % Generator
    \node[draw, rectangle, fill=laranjacel!30, text width=2cm, align=center] (gen) at (0,2) {Generator};

    % Discriminator
    \node[draw, rectangle, fill=azulclaro!30, text width=2cm, align=center] (disc) at (0,0) {Discriminator};

    % Real data
    \node[draw, ellipse, fill=verdeacido!30] (real) at (-3,0) {Real};

    % Fake data
    \draw[->, thick] (gen) -- node[right, font=\tiny] {Fake} (disc);
    \draw[->, thick] (real) -- (disc);

    % Noise
    \node[above] at (0,3) {Noise $z$};
    \draw[->, thick] (0,2.8) -- (gen);

    % Loss
    \draw[->, thick, red] (disc) -- node[right, font=\tiny] {Loss} ++(2,-1);
    \draw[->, thick, red, bend right] (2,-1) -- (gen);
\end{tikzpicture}
\end{center}

\begin{alertblock}{Exemplo}
MegaMolBART (transformers) gera moléculas drug-like otimizadas
\end{alertblock}
\end{columns}
\end{frame}

\begin{frame}{IA Interpretável e Explicável}
\importante{
Modelos de caixa-preta são problemáticos em biologia e medicina!
}

\vspace{0.3cm}

\begin{block}{Abordagens para Interpretabilidade}
\begin{enumerate}
    \item \textbf{Feature importance:} quais genes são mais importantes?
    \item \textbf{SHAP values:} contribuição de cada feature
    \item \textbf{Attention maps:} o que o modelo está "olhando"
    \item \textbf{Layer-wise relevance propagation (LRP)}
    \item \textbf{Modelos híbridos:} integrar conhecimento biológico
\end{enumerate}
\end{block}

\vspace{0.3cm}

\begin{exampleblock}{Inferência Causal vs Correlação}
\begin{itemize}
    \item Correlação: $A$ e $B$ variam juntos
    \item Causação: $A \rightarrow B$ (mudança em $A$ causa mudança em $B$)
    \item Ferramentas: causal graphs, do-calculus, experimentos perturbacionais
\end{itemize}
\end{exampleblock}
\end{frame}

% ============================================================================
% SEÇÃO 5: DOENÇAS COMPLEXAS E ENVELHECIMENTO
% ============================================================================
\section{Doenças Complexas e Envelhecimento}

\begin{frame}{15.4 Doenças Complexas como Perturbações de Rede}
\begin{block}{Abordagem de Network Medicine}
Doenças complexas (câncer, diabetes, Alzheimer) não resultam de defeito em gene único, mas de perturbações em redes moleculares.
\end{block}

\begin{center}
\begin{tikzpicture}[scale=0.7]
    % Healthy network
    \begin{scope}
        \foreach \i in {1,...,8} {
            \pgfmathsetmacro{\angle}{360/8 * \i}
            \node[draw, circle, fill=verdeacido!30, minimum size=0.5cm] (h\i) at (\angle:1.5) {};
        }
        \foreach \i in {1,...,8} {
            \pgfmathsetmacro{\j}{mod(\i, 8) + 1}
            \draw[thick] (h\i) -- (h\j);
        }
        \node[below] at (0,-2.5) {Saudável};
    \end{scope}

    % Disease network
    \begin{scope}[xshift=5cm]
        \foreach \i in {1,...,8} {
            \pgfmathsetmacro{\angle}{360/8 * \i}
            \pgfmathsetmacro{\col}{(\i == 2 || \i == 3) ? "red" : "azulclaro"}
            \node[draw, circle, fill=\col!30, minimum size=0.5cm] (d\i) at (\angle:1.5) {};
        }
        \foreach \i in {1,...,8} {
            \pgfmathsetmacro{\j}{mod(\i, 8) + 1}
            \pgfmathsetmacro{\col}{(\i == 2 || \i == 3) ? "red" : "black"}
            \draw[thick, \col] (d\i) -- (d\j);
        }
        \node[below] at (0,-2.5) {Doença};
    \end{scope}
\end{tikzpicture}
\end{center}

\textbf{Nós vermelhos:} genes mutados/desregulados\\
\textbf{Propagação:} perturbação se espalha pela rede
\end{frame}

\begin{frame}{Multi-Omics para Envelhecimento}
\begin{columns}
\column{0.5\textwidth}
\textbf{Hallmarks of Aging:}
\begin{enumerate}\small
    \item Instabilidade genômica
    \item Atrito de telômeros
    \item Alterações epigenéticas
    \item Perda de proteostase
    \item Senescência celular
    \item Disfunção mitocondrial
    \item Desregulação nutricional
    \item Exaustão de células-tronco
    \item Comunicação intercelular alterada
\end{enumerate}

\column{0.5\textwidth}
\textbf{Abordagem Multi-Omic:}
\begin{itemize}
    \item Genômica (mutações)
    \item Epigenômica (metilação)
    \item Transcriptômica (expressão)
    \item Proteômica (proteostase)
    \item Metabolômica (NAD$^+$, etc.)
\end{itemize}

\vspace{0.3cm}

\begin{alertblock}{Relógios Epigenéticos}
DNA methylation age: preditor de idade biológica (Horvath clock)
\end{alertblock}
\end{columns}
\end{frame}

\begin{frame}{Vias de Longevidade}
\begin{block}{Pathways Conservados}
Vias moleculares que modulam longevidade em múltiplas espécies:
\end{block}

\begin{center}
\begin{tikzpicture}[node distance=2.5cm, auto,
    box/.style={rectangle, draw, text width=2.2cm, text centered, rounded corners, minimum height=0.8cm, font=\small}]

    \node[box, fill=roxodna!30] (igf) {IGF-1/Insulina};
    \node[box, fill=laranjacel!30, right of=igf] (tor) {mTOR};
    \node[box, fill=verdeacido!30, right of=tor] (sirt) {Sirtuínas};
    \node[box, fill=azulclaro!30, below of=tor] (ampk) {AMPK};

    \draw[->, thick] (igf) -- (tor);
    \draw[->, thick] (tor) -- (sirt);
    \draw[->, thick] (ampk) -- (tor);
    \draw[-|, thick, red] (sirt) to[bend left] (igf);

    \node[below=0.5cm of ampk, font=\small] {Longevidade $\uparrow$};
\end{tikzpicture}
\end{center}

\textbf{Intervenções:}
\begin{itemize}
    \item Restrição calórica
    \item Rapamicina (inibe mTOR)
    \item Metformina (ativa AMPK)
    \item NAD$^+$ boosters (ativam sirtuínas)
    \item Senolytics (eliminam células senescentes)
\end{itemize}
\end{frame}

\begin{frame}{Redes de Comorbidade}
\begin{block}{Comorbidity Networks}
Grafos representando co-ocorrência de doenças em populações, revelando mecanismos compartilhados.
\end{block}

\begin{center}
\begin{tikzpicture}[scale=0.8]
    % Diseases as nodes
    \node[draw, circle, fill=red!30, minimum size=1cm] (dm) at (0,0) {DM2};
    \node[draw, circle, fill=orange!30, minimum size=0.8cm] (htn) at (2.5,1) {HTN};
    \node[draw, circle, fill=yellow!30, minimum size=0.7cm] (cad) at (2.5,-1) {CAD};
    \node[draw, circle, fill=green!30, minimum size=0.6cm] (ckd) at (-2,1) {CKD};
    \node[draw, circle, fill=blue!30, minimum size=0.5cm] (stroke) at (-2,-1) {AVC};

    % Edges (comorbidity)
    \draw[thick, line width=3pt] (dm) -- (htn);
    \draw[thick, line width=2.5pt] (dm) -- (cad);
    \draw[thick, line width=2pt] (dm) -- (ckd);
    \draw[thick, line width=1.5pt] (htn) -- (cad);
    \draw[thick, line width=1.5pt] (htn) -- (stroke);
    \draw[thick, line width=1pt] (ckd) -- (stroke);
\end{tikzpicture}
\end{center}

\textbf{Insight:} diabetes mellitus tipo 2 (DM2) é hub central\\
\textbf{Aplicação:} priorizar intervenções, predizer risco
\end{frame}

\begin{frame}{Sistemas para Preparação Pandêmica}
\importante{
COVID-19 demonstrou necessidade de biologia de sistemas para resposta rápida!
}

\vspace{0.3cm}

\begin{block}{Abordagens Sistêmicas}
\begin{enumerate}
    \item \textbf{Vigilância genômica:} sequenciamento em tempo real de variantes
    \item \textbf{Modelos epidemiológicos:} previsão de disseminação (SIR, SEIR)
    \item \textbf{Multi-omics de hospedeiro:} identificar fatores de risco
    \item \textbf{Drug repurposing:} IA para identificar drogas existentes
    \item \textbf{Design de vacinas:} epitopos computacionais
    \item \textbf{Network medicine:} entender síndrome pós-COVID
\end{enumerate}
\end{block}

\vspace{0.3cm}

\textbf{Lição:} integração de dados heterogêneos + modelos + infraestrutura computacional
\end{frame>

\begin{frame}{Medicina de Precisão: Futuro}
\begin{columns}
\column{0.5\textwidth}
\textbf{De População para Indivíduo:}
\begin{itemize}
    \item Genoma pessoal (WGS)
    \item Multi-omics longitudinal
    \item Microbioma
    \item Wearables (monitoramento)
    \item Histórico médico completo
\end{itemize}

\vspace{0.3cm}

\textbf{Desafios:}
\begin{itemize}
    \item Custo
    \item Interpretação clínica
    \item Privacidade
    \item Equidade de acesso
\end{itemize}

\column{0.5\textwidth}
\begin{center}
\begin{tikzpicture}[scale=0.7]
    % Individual
    \node[draw, circle, fill=azulclaro!30, minimum size=1.5cm] (person) at (0,0) {Paciente};

    % Data sources
    \node[draw, ellipse, fill=roxodna!30, font=\tiny] (genome) at (-2,2) {Genoma};
    \node[draw, ellipse, fill=laranjacel!30, font=\tiny] (omics) at (0,2.5) {Multi-\\omics};
    \node[draw, ellipse, fill=verdeacido!30, font=\tiny] (micro) at (2,2) {Micro-\\bioma};
    \node[draw, ellipse, fill=yellow!30, font=\tiny] (wear) at (2.5,0) {Wear-\\ables};
    \node[draw, ellipse, fill=red!30, font=\tiny] (ehr) at (-2.5,0) {EHR};

    % Connections
    \foreach \src in {genome, omics, micro, wear, ehr} {
        \draw[->, thick] (\src) -- (person);
    }

    % Output
    \draw[->, ultra thick, blue] (person) -- ++(0,-2) node[below, font=\small] {Tratamento\\Personalizado};
\end{tikzpicture}
\end{center}
\end{columns}
\end{frame>

% ============================================================================
% SEÇÃO 6: DIREÇÕES FUTURAS
% ============================================================================
\section{Direções Futuras}

\begin{frame}{15.5 Direções Futuras}
\begin{block}{O Futuro da Biologia de Sistemas}
A área está em rápida evolução. Várias tecnologias e abordagens prometem revolucionar nossa compreensão da vida.
\end{block}

\begin{columns}
\column{0.5\textwidth}
\textbf{Tecnologias Emergentes:}
\begin{itemize}
    \item Laboratórios automatizados
    \item Plataformas robóticas
    \item Microfluídica
    \item Organoides e organs-on-chip
    \item Edição genômica (CRISPR)
    \item Biologia sintética
\end{itemize>

\column{0.5\textwidth}
\textbf{Computação:}
\begin{itemize}
    \item Computação quântica
    \item Edge computing
    \item Federated learning
    \item Blockchain para dados
    \item Cloud-based pipelines
\end{itemize}
\end{columns}
\end{frame}

\begin{frame}{Laboratórios Automatizados e Robótica}
\begin{center}
\begin{tikzpicture}[scale=0.8]
    % Self-driving lab
    \draw[thick] (0,0) rectangle (8,4);
    \node[above] at (4,4) {\textbf{Laboratório Autônomo}};

    % Components
    \node[draw, rectangle, fill=roxodna!30, text width=1.5cm, align=center, font=\tiny] at (1,3) {Robô\\Líquidos};
    \node[draw, rectangle, fill=laranjacel!30, text width=1.5cm, align=center, font=\tiny] at (4,3) {Sequen-\\ciador};
    \node[draw, rectangle, fill=verdeacido!30, text width=1.5cm, align=center, font=\tiny] at (7,3) {Microscópio\\Automatizado};

    % AI controller
    \node[draw, rectangle, fill=azulclaro!30, text width=2cm, align=center] at (4,1) {Controlador\\IA};

    % Arrows
    \draw[->, thick] (4,1.5) -- (1,2.5);
    \draw[->, thick] (4,1.5) -- (4,2.5);
    \draw[->, thick] (4,1.5) -- (7,2.5);

    % Feedback
    \draw[->, thick, red, bend right=45] (7,2.5) to node[right, font=\tiny] {Dados} (4,0.5);
\end{tikzpicture}
\end{center>

\textbf{Vantagens:}
\begin{itemize}
    \item Experimentos 24/7
    \item Alta reprodutibilidade
    \item Exploração massiva do espaço de parâmetros
    \item Loop fechado: IA propõe experimentos $\rightarrow$ robô executa $\rightarrow$ análise $\rightarrow$ nova hipótese
\end{itemize}

\textbf{Exemplos:} Emerald Cloud Lab, Strateos, Ada (Carnegie Mellon)
\end{frame}

\begin{frame}{Geração de Hipóteses por IA}
\begin{block}{IA Criativa}
Sistemas de IA que não apenas analisam dados, mas propõem novas hipóteses biológicas testáveis.
\end{block}

\vspace{0.3cm}

\begin{exampleblock}{Abordagens}
\begin{enumerate}
    \item \textbf{Mineração de literatura:} extrair relações de milhões de papers
    \item \textbf{Knowledge graphs:} raciocínio sobre entidades e relações
    \item \textbf{Reinforcement learning:} explorar espaço de experimentos
    \item \textbf{Generative models:} propor novas moléculas, vias
    \item \textbf{Causal inference:} identificar relações causais
\end{enumerate}
\end{exampleblock}

\vspace{0.3cm}

\importante{
Humanos continuam essenciais: interpretar, validar, contextualizar!
}
\end{frame}

\begin{frame}{Digital Twins para Medicina Personalizada}
\definicao{Digital Twin}{
Réplica computacional dinâmica de um paciente individual, integrando dados multi-omics, fisiológicos e clínicos em tempo real.
}

\vspace{0.3cm}

\begin{center}
\begin{tikzpicture}[scale=0.8]
    % Real patient
    \node[draw, rectangle, fill=azulclaro!30, text width=2cm, align=center] (patient) at (0,0) {Paciente\\Real};

    % Digital twin
    \node[draw, rectangle, fill=laranjacel!30, text width=2cm, align=center] (twin) at (5,0) {Digital\\Twin};

    % Data flow
    \draw[->, thick, blue] (patient) -- node[above, font=\small] {Dados} (twin);

    % Simulation
    \node[draw, cloud, fill=verdeacido!30, text width=1.5cm, align=center, font=\small] (sim) at (5,2.5) {Simulações};
    \draw[->, thick] (twin) -- (sim);

    % Predictions
    \draw[->, thick, red] (sim) to[bend left] node[right, font=\small] {Predições} (patient);

    % Wearables
    \node[font=\tiny] at (0,-1.5) {Wearables, EHR,\\omics, microbioma};
\end{tikzpicture}
\end{center>

\textbf{Aplicações:}
\begin{itemize}
    \item Testar tratamentos \textit{in silico}
    \item Predizer resposta a drogas
    \item Otimizar dosagem
    \item Antecipar eventos adversos
\end{itemize}
\end{frame}

\begin{frame}{Computação Quântica para Simulação Biológica}
\begin{block}{Computação Quântica}
Computadores quânticos exploram superposição e entrelaçamento para resolver problemas intratáveis classicamente.
\end{block}

\textbf{Aplicações Potenciais em Biologia:}
\begin{enumerate}
    \item \textbf{Simulação quântica de moléculas:} dinâmica eletrônica precisa
    \item \textbf{Folding de proteínas:} explorar espaço conformacional exponencial
    \item \textbf{Drug discovery:} otimização combinatorial
    \item \textbf{Análise de grafos:} redes biológicas massivas
    \item \textbf{Machine learning quântico:} aceleração de algoritmos
\end{enumerate}

\vspace{0.3cm}

\begin{alertblock}{Realidade Atual}
Ainda estamos na era NISQ (Noisy Intermediate-Scale Quantum). Aplicações práticas em biologia ainda distantes, mas pesquisa ativa!
\end{alertblock}
\end{frame}

\begin{frame}{Integração de Microbioma e Fatores Ambientais}
\importante{
O hospedeiro não pode ser estudado isoladamente: microbioma e ambiente são partes integrantes do sistema!
}

\vspace{0.3cm}

\begin{columns}
\column{0.5\textwidth}
\textbf{Microbioma:}
\begin{itemize}
    \item 10$\times$ mais células que humanas
    \item Metabolismo complementar
    \item Modulação imune
    \item Eixo intestino-cérebro
    \item Resistência a patógenos
\end{itemize}

\vspace{0.3cm}

\textbf{Metagenômica:}
\begin{itemize}
    \item 16S rRNA-seq (taxonomia)
    \item Shotgun metagenomics (função)
    \item Metatranscriptomics
\end{itemize}

\column{0.5\textwidth}
\textbf{Exposoma:}
\begin{itemize}
    \item Dieta
    \item Poluentes
    \item Estresse
    \item Atividade física
    \item Sono
    \item Medicamentos
\end{itemize}

\vspace{0.3cm}

\begin{exampleblock}{Holobionte}
Organismo + microbioma = unidade evolutiva e funcional
\end{exampleblock}
\end{columns}
\end{frame>

\begin{frame}{Ciência Aberta e Compartilhamento de Dados}
\begin{block}{Movimento Open Science}
Transparência, reprodutibilidade e acesso aberto são cruciais para acelerar descobertas.
\end{block}

\textbf{Princípios FAIR:}
\begin{itemize}
    \item \textbf{F}indable: dados encontráveis
    \item \textbf{A}ccessible: acessíveis
    \item \textbf{I}nteroperable: interoperáveis
    \item \textbf{R}eusable: reutilizáveis
\end{itemize}

\vspace{0.3cm}

\textbf{Recursos:}
\begin{itemize}
    \item Repositórios públicos (GEO, SRA, PDB, etc.)
    \item Pré-prints (bioRxiv, medRxiv)
    \item Código aberto (GitHub)
    \item Notebooks reprodutíveis (Jupyter)
    \item Protocolos detalhados (protocols.io)
\end{itemize}

\importante{
Colaboração global é essencial para enfrentar desafios complexos!
}
\end{frame}

\begin{frame}{Considerações Éticas}
\begin{alertblock}{Questões Éticas Emergentes}
\begin{enumerate}
    \item \textbf{Privacidade:} genomas pessoais, re-identificação
    \item \textbf{Consentimento informado:} uso secundário de dados
    \item \textbf{Equidade:} acesso a tecnologias avançadas (precision medicine gap)
    \item \textbf{Viés algorítmico:} datasets não-representativos
    \item \textbf{Edição germline:} CRISPR em embriões
    \item \textbf{Biologia sintética:} biossegurança, dual use
    \item \textbf{Propriedade intelectual:} patentes de genes, dados
\end{enumerate}
\end{alertblock}

\vspace{0.3cm}

\importante{
Cientistas têm responsabilidade de considerar implicações sociais de seu trabalho!
}
\end{frame}

% ============================================================================
% SEÇÃO 7: EXERCÍCIOS
% ============================================================================
\section{Exercícios}

\begin{frame}{Exercícios - Parte 1}
\begin{block}{Questão 1: Modelagem Multi-Escala}
Descreva como você modelaria o metabolismo da glicose integrando os níveis:
\begin{enumerate}
    \item Molecular (cinética de hexoquinase)
    \item Celular (hepatócito)
    \item Órgão (fígado)
\end{enumerate}
Quais são os principais desafios de acoplamento entre escalas?
\end{block}

\begin{block}{Questão 2: Single-Cell RNA-seq}
Dado um dataset scRNA-seq de um tumor:
\begin{enumerate}
    \item Como você identificaria diferentes tipos celulares?
    \item Como inferiria comunicação entre células tumorais e células imunes?
    \item Qual a importância da informação espacial?
\end{enumerate}
\end{block}
\end{frame}

\begin{frame}{Exercícios - Parte 2}
\begin{block}{Questão 3: Machine Learning}
\begin{enumerate}
    \item Quando usar modelos de ML vs. modelos mecanísticos?
    \item Explique o problema de interpretabilidade em redes neurais profundas.
    \item Como você validaria um modelo de ML para predição clínica?
\end{enumerate}
\end{block}

\begin{alertblock}{Projeto Prático}
Escolha um tópico emergente (célula única, IA, envelhecimento) e:
\begin{enumerate}
    \item Revise a literatura recente (últimos 2 anos)
    \item Identifique um dataset público relevante
    \item Aplique pelo menos uma técnica computacional
    \item Gere visualizações e interprete resultados
    \item Discuta limitações e direções futuras
\end{enumerate>
\end{alertblock}
\end{frame}

% ============================================================================
% SEÇÃO 8: REFERÊNCIAS
% ============================================================================
\section{Referências}

\begin{frame}{Referências Principais}
\begin{thebibliography}{99}
\bibitem{jumper2021} Jumper J. et al. (2021). Highly accurate protein structure prediction with AlphaFold. \textit{Nature}.

\bibitem{wolf2018} Wolf F.A. et al. (2018). SCANPY: large-scale single-cell gene expression data analysis. \textit{Genome Biology}.

\bibitem{lopez2020} López-Otín C. et al. (2023). Hallmarks of aging: An expanding universe. \textit{Cell}.

\bibitem{barabasi2011} Barabási A.L. et al. (2011). Network medicine: a network-based approach to human disease. \textit{Nature Reviews Genetics}.

\bibitem{moses2021} Moses C. et al. (2021). Museum of spatial transcriptomics. \textit{Nature Methods}.
\end{thebibliography}
\end{frame}

\begin{frame}{Leitura Complementar}
\begin{block}{Artigos de Revisão Recentes}
\begin{itemize}
    \item Integrative multi-omics approaches (Nature Reviews, 2023)
    \item Single-cell multi-omics technologies (Nature Biotechnology, 2023)
    \item AI for drug discovery (Nature Machine Intelligence, 2024)
    \item Digital twins in healthcare (Nature Medicine, 2023)
    \item Longevity and healthspan interventions (Cell, 2023)
\end{itemize}
\end{block>

\begin{block}{Recursos Online}
\begin{itemize}
    \item \textbf{AlphaFold DB:} \url{https://alphafold.ebi.ac.uk}
    \item \textbf{Single Cell Portal:} \url{https://singlecell.broadinstitute.org}
    \item \textbf{Human Cell Atlas:} \url{https://www.humancellatlas.org}
    \item \textbf{PyTorch Geometric:} \url{https://pytorch-geometric.readthedocs.io}
\end{itemize}
\end{block}
\end{frame}

% ============================================================================
% SLIDE FINAL: CONCLUSÃO DO CURSO
% ============================================================================
\begin{frame}
\begin{center}
{\Huge Obrigado!}

\vspace{0.5cm}

{\Large Capítulo 15: Emerging Topics}

\vspace{1cm}

\begin{block}{Conclusão do Curso}
Este curso cobriu desde fundamentos de sistemas biológicos até os tópicos mais avançados e emergentes da biologia de sistemas. Vocês agora têm as ferramentas conceituais e computacionais para:
\begin{itemize}
    \item Modelar sistemas biológicos em múltiplos níveis
    \item Analisar dados omics em larga escala
    \item Integrar conhecimento multi-escala e multi-omic
    \item Aplicar IA e aprendizado de máquina
    \item Contribuir para o futuro da medicina de precisão
\end{itemize}
\end{block}

\vspace{0.5cm}

\textit{O futuro da biologia de sistemas é brilhante.\\Boa sorte em suas jornadas científicas!}
\end{center}
\end{frame}

\end{document}
